\par\Needspace{4\baselineskip}
\begin{center}
\hypertarget{native-052-article035}{}\pdfbookmark[2]{35.}{article-035}35.
\end{center}
Unser \textsc{Bessel} leitete die analytischen Ausdrücke für die ersten drei Koeffizienten der zweiten Reihe, $w^{\prime}$, $w^{\prime\prime}$, $w^{\prime\prime\prime}$, her und fügte für die Zahlenwerte der ersten beiden, $w^{\prime}$ und $w^{\prime\prime}$, eine Tafel für jeden Grad des Arguments $w$ hinzu (\textsc{von Zach} Monatliche Correspondenz, Bd. XII, S. 197). Für den ersten Koeffizienten $w^{\prime}$ gab es bereits eine von \textsc{Simpson} berechnete Tafel, die dem oben angeführten Werk des berühmten \textsc{Olbers} beigefügt ist. In den meisten Fällen kann man mit diesem Verfahren mithilfe der \textsc{Bessel}schen Tafel die wahre Anomalie aus der Zeit mit hinreichender Genauigkeit \SourceLeaf{053}bestimmen. Was noch zu wünschen übrig bleibt, lässt sich im Wesentlichen auf die folgenden Punkte zurückführen.

I. Beim inversen Problem, nämlich der Bestimmung der Zeit aus der wahren Anomalie, muss man zu einem gewissermaßen indirekten Verfahren greifen und $w$ aus $v$ durch Probieren herleiten. Um diese Schwierigkeit zu beheben, müsste man die erste Reihe ebenso behandeln wie die zweite. Man sieht leicht, dass $-v^{\prime}$ dieselbe Funktion von $v$ ist wie $w^{\prime}$ von $w$, sodass man die Tafel für $w^{\prime}$ allein durch Umkehrung des Vorzeichens auch für $v^{\prime}$ verwenden kann. Es wäre also nur noch eine Tafel für $v^{\prime\prime}$ erforderlich, um beide Probleme mit gleicher Genauigkeit lösen zu können.

II. Es können durchaus Fälle vorkommen, in denen die Exzentrizität so wenig von eins abweicht, dass die oben erläuterten allgemeinen Verfahren keine hinreichende Genauigkeit zu liefern scheinen, die Abweichung aber noch zu groß ist, als dass man beim eben skizzierten besonderen Verfahren den Einfluss der dritten und höheren Potenzen von $\delta$ unbedenklich vernachlässigen könnte. Insbesondere bei der hyperbolischen Bewegung sind Fälle möglich, in denen ein Fehler von mehreren Bogensekunden unvermeidlich ist, wenn man nur die üblichen auf sieben Stellen berechneten Tafeln verwendet, gleichviel, ob man nach den allgemeinen Verfahren oder nach diesem besonderen Verfahren rechnet. Auch wenn solche Fälle in der Praxis selten vorkommen, könnte man doch mit Recht eine Lücke darin sehen, wenn sich die wahre Anomalie nur mithilfe größerer Tafeln in \emph{allen} Fällen bis auf $0^{\prime\prime}{,}1$ oder wenigstens $0^{\prime\prime}{,}2$ bestimmen ließe, denn diese Tafeln gehören bekanntlich zu den selteneren Büchern. Wir hoffen daher, dass es sich lohnt, ein besonderes Verfahren darzulegen, das wir schon seit Langem anwenden. Ein weiterer Vorzug dieses Verfahrens ist seine allgemeine Anwendbarkeit, zumindest was die Abweichung der Exzentrizität von eins angeht.

\par\Needspace{4\baselineskip}
\begin{center}
\hypertarget{native-053-article036}{}\pdfbookmark[2]{36.}{article-036}36.
\end{center}
Bevor wir dieses Verfahren erläutern, sei bemerkt, dass die Unsicherheit der oben angegebenen allgemeinen Verfahren bei Bahnen, die sich einer Parabel nähern, von selbst verschwindet, sobald $E$ oder $F$ eine beträchtliche Größe erreicht. Das geschieht erst in großen Entfernungen von der Sonne. Um dies zu zeigen, schreiben wir den größtmöglichen Fehler in der Ellipse, den wir in Art. 32, IV als $\frac{3\omega ea\sin v}{\lambda r}\cdot206265^{\prime\prime}$ gefunden haben, in der Form
\[
\frac{3\omega e\BareRoot{(1-ee)}.\sin E}{\lambda(1-e\cos E)^2}\cdot206265^{\prime\prime},
\]
\SourceLeaf{054}Daraus geht unmittelbar hervor, dass sich der Fehler stets in engen Grenzen hält, sobald $E$ einen beträchtlichen Wert erreicht hat oder, anders gesagt, sobald $\cos E$ weiter von eins abweicht, wie groß die Exzentrizität auch sein mag. Noch deutlicher zeigt dies die folgende Tafel, in der wir den größten Zahlenwert dieser Formel für einige feste Werte berechnet haben (bei sieben Dezimalstellen).
\begin{center}
\begin{tabular}{r|lr}
$E=10^0$ & größter Fehler $=$ & $3^{\prime\prime}{,}04$\\
$20$ & & $0{,}76$\\
$30$ & & $0{,}34$\\
$40$ & & $0{,}19$\\
$50$ & & $0{,}12$\\
$60$ & & $0{,}08$
\end{tabular}
\end{center}
\noindent Ähnlich verhält es sich bei der Hyperbel, wie unmittelbar zu erkennen ist, wenn man den in Art. 32, VII hergeleiteten Ausdruck in der Form $\frac{\omega\cos F.(\cos F+3e\sin F)\BareRoot{(ee-1)}}{\lambda(e-\cos F)^2}206265^{\prime\prime}$ schreibt. Die folgende Tafel gibt die größten Werte dieses Ausdrucks für einige feste Werte von $F$ an.
\begin{center}
\begin{tabular}{r|r|r|r}
\hline
$F$ & \multicolumn{2}{c|}{$u$} & größter Fehler\\
\hline
$10^0$ & $1{,}192$ & $0{,}839$ & $8^{\prime\prime}{,}66$\\
$20$ & $1{,}428$ & $0{,}700$ & $1{,}38$\\
$30$ & $1{,}732$ & $0{,}577$ & $0{,}47$\\
$40$ & $2{,}144$ & $0{,}466$ & $0{,}22$\\
$50$ & $2{,}747$ & $0{,}364$ & $0{,}11$\\
$60$ & $3{,}732$ & $0{,}268$ & $0{,}06$\\
$70$ & $5{,}671$ & $0{,}176$ & $0{,}03$
\end{tabular}
\end{center}
\noindent Überschreitet also $E$ oder $F$ den Wert $40^0$ oder $50^0$ (was allerdings bei Bahnen, die sich nur wenig von einer Parabel unterscheiden, kaum vorkommen wird, da sich Himmelskörper auf solchen Bahnen in so großen Entfernungen von der Sonne meist unseren Blicken entziehen), so gibt es keinen Grund, das allgemeine Verfahren aufzugeben. Übrigens würden in einem solchen Fall auch die in Art. 34 behandelten Reihen zu langsam konvergieren. Dass das nun zu erläuternde Verfahren vor allem für Fälle geeignet ist, in denen $E$ oder $F$ noch nicht über mäßige Werte hinausgewachsen ist, kann daher keineswegs als Mangel gelten.
